\section{Resultados: valor da seleção por perfil}
\label{sec:ba2026}

\subsection{Quais candidatas oferecem compromissos competitivos?}

A Tabela~\ref{tab:ba2026_portfolio} resume os 290 casos série--origem. Sete casos não têm demanda no teste: um na primeira origem e seis na segunda. Seus custos permanecem na média dos 290 casos; o atendimento médio utiliza os 283 casos com denominador positivo. O atendimento ponderado é a razão entre a soma das unidades atendidas e a soma da demanda nos blocos avaliados, após a média das sementes. As origens repetidas não ampliam a amostra para 290 séries independentes.

\begin{table}[htbp]
\centering\small
\caption{Bahia, protocolo temporal: médias no teste de sete ciclos. Custo em unidades do cenário, calculado nos 290 casos. Atendimento médio e contagem abaixo da meta consideram os 283 casos com demanda positiva; sete casos têm atendimento indefinido.}
\label{tab:ba2026_portfolio}
\begin{tabular}{@{}lrrrr@{}}\toprule
Política & Custo & Atendimento & Ponderado & Abaixo de 0,70\\\midrule
EOQCycle &166,77&0,7787&0,6265&101\\
sSCycle &164,23&0,8037&0,6802&90\\
CappedBaseStock &218,11&0,6631&0,4948&148\\
GA &143,85&0,6953&0,5457&140\\
PPO &134,23&0,5514&0,4189&171\\
GA--DQN &145,75&0,6802&0,5248&144\\
DQN (controle) &164,88&0,5718&0,4033&164\\
\bottomrule\end{tabular}
\end{table}

PPO tem o menor custo descritivo, acompanhado do menor atendimento médio. Apenas EOQCycle e sSCycle superam 0,70 no conjunto dos testes; nenhuma política alcança essa meta no atendimento ponderado. Entre essas duas referências, sSCycle apresenta menor custo e maior atendimento médio. Essa comparação é um diagnóstico do teste, não uma escolha retrospectiva a ser atribuída ao seletor.

GA--DQN custa 11,60\% menos que DQN e aumenta o atendimento médio de 0,5718 para 0,6802. Frente ao GA isolado, porém, custa 1,32\% mais e atende menos, com média de 0,6953 para o GA. O refinamento híbrido não oferece vantagem agregada nessas duas dimensões sobre seu componente evolutivo. A diferença varia entre origens: na segunda, seu custo é ligeiramente menor que o do GA, também com menor atendimento. A comparação descreve o comportamento das implementações avaliadas, sem isolar causalmente o benefício das demonstrações ou do gatilho.

\begin{figure}[htbp]
\centering\includegraphics[width=\linewidth]{figures/custo_servico.pdf}
\caption{Compromisso custo--serviço das sete candidatas na Bahia. Médias de teste em 290 casos série--origem para custo e 283 com demanda positiva para atendimento. A fronteira é descritiva: não utiliza o teste para escolher a recomendação nem representa dominância estatística ou por série.}
\label{fig:cost_service}
\end{figure}

A Figura~\ref{fig:cost_service} distingue economia aparente por baixo atendimento de um compromisso competitivo. A restrição usada pelo seletor continua sendo avaliada na validação, sem recorrer a essa fronteira de teste.

\subsection{O perfil melhora a decisão frente à seleção global?}

\begin{table}[htbp]
\centering\small
\caption{Escolha na validação das demais lojas e avaliação no teste da loja reservada, com meta 0,70. Há 145 casos por linha; atendimento definido em 144 casos na origem 17 e 139 na origem 23.}
\label{tab:ba2026_selection}
\begin{tabular}{@{}llrrr@{}}\toprule
Treino & Estratégia & Custo & Atendimento & Abaixo de 0,70\\\midrule
17 ciclos & Global &180,58&0,5711&101\\
17 ciclos & Perfil &186,79&0,6475&82\\
23 ciclos & Global &142,03&0,9073&18\\
23 ciclos & Perfil &142,03&0,9073&18\\
\bottomrule\end{tabular}
\end{table}

Na primeira origem, a regra global escolhe GA para todas as séries. Por perfil, as escolhas são GA em 65 casos, sSCycle em 75 e EOQCycle em cinco. O maior atendimento custa mais, e ambas as estratégias ficam abaixo da meta no teste. Na segunda origem, os dois seletores escolhem sSCycle em todos os casos. A Tabela~\ref{tab:ba2026_selection} evidencia, portanto, tanto mudança temporal da escolha global quanto ausência de diferenciação contextual na segunda janela.

No conjunto das origens, o custo médio passa de 161,31 na estratégia global para 164,41 por perfil: aumento de 1,93\%. O atendimento médio cresce de 0,7362 para 0,7751, diferença de 3,89 pontos percentuais; o ponderado passa de 0,5750 para 0,6395. Os casos abaixo da meta diminuem de 119 para 100, entre os 283 com atendimento definido. A estimativa pontual combina custo maior e melhor serviço, com a diferença concentrada na primeira origem. Essa melhora não assegura atendimento mínimo individual nem equivale a cumprir uma meta ponderada de 0,70.

Uma reamostragem exploratória pareada de 115 lojas, com 5.000 sorteios e todas as origens e produtos de cada loja preservados, produz intervalo condicional de 95\% de $[-0{,}32;4{,}16]\%$ para o aumento de custo por perfil. Os controladores e seletores permanecem fixos nos sorteios. Para a diferença de atendimento, o mesmo procedimento produz $[2{,}78;5{,}07]$ pontos percentuais. O intervalo de custo inclui zero; o de atendimento permanece positivo. Ambos são condicionais e não representam incerteza de novos treinamentos ou de novas trajetórias de demanda; dependência entre lojas que compartilham produtos também permanece possível.

\subsection{Sensibilidade ao limiar de escolha}

A Tabela~\ref{tab:threshold_sensitivity} apresenta todos os limiares avaliados. Em 0,80, a estratégia por perfil tem atendimento médio ligeiramente menor e custo maior que a global. Em 0,85, ambas selecionam sSCycle nos 290 casos, embora a frequência de contingência seja diferente: uma escolha pode cumprir o critério de seu grupo e não cumprir o critério agregado, ou o contrário. Em 0,90 e 0,95, ambas usam contingências em todos os casos. Portanto, a diferenciação contextual não preserva a vantagem de atendimento em todas as metas.

\begin{table}[htbp]\centering\small
\caption{Sensibilidade da seleção, sem retreinamento. Custo em 290 casos; atendimento médio em 283. Contingências indicam ausência de candidata viável na validação; não informam cumprimento de serviço no teste. G: global; P: por perfil.}
\label{tab:threshold_sensitivity}
\begin{tabular}{@{}rrrrrr@{}}\toprule
Meta & Custo G & Custo P & Atend. G & Atend. P & Conting. G/P\\\midrule
0,70 &161,31&164,41&0,7362&0,7751&0/0\\
0,80 &164,33&164,53&0,8037&0,8025&145/223\\
0,85 &164,23&164,23&0,8037&0,8037&290/225\\
0,90 &164,23&164,23&0,8037&0,8037&290/290\\
0,95 &164,23&164,23&0,8037&0,8037&290/290\\
\bottomrule\end{tabular}
\end{table}

As candidatas continuam calibradas com referência de 0,70. A inexistência de uma opção viável em uma meta maior não demonstra impossibilidade de atendê-la com outro ajuste ou portfólio. A sensibilidade examina a regra de seleção sobre a matriz disponível.

\begin{figure}[htbp]
\centering\includegraphics[width=\linewidth]{figures/distribuicao_atendimento.pdf}
\caption{Distribuição do atendimento das estratégias no teste, por origem. Cada caixa resume séries com demanda positiva; a linha tracejada indica 0,70. Há 144 casos por estratégia na primeira origem e 139 na segunda. As mesmas séries participam das duas estratégias; observações e origens não são independentes.}
\label{fig:service_distribution}
\end{figure}

A Figura~\ref{fig:service_distribution} mostra por que a comparação não se esgota na média: a primeira janela concentra violações de serviço, enquanto a segunda tem distribuição e escolha coincidentes entre estratégias.

\subsection{As escolhas e o atendimento são estáveis entre períodos?}

Os perfis são recalculados no prefixo disponível em cada origem. Com 17 ciclos, há 65 séries Sparse High Impact e 80 Unstable Trend; com 23, são 60, 69 e 16 High Volatility Seasonal. Essa mudança reflete a janela e as regras de caracterização, inclusive a indisponibilidade da autocorrelação na primeira origem. Não demonstra mudança de regime econômico em todas essas séries.


\begin{figure}[htbp]
\centering\includegraphics[width=\linewidth]{figures/escolhas_perfil.pdf}
\caption{Políticas escolhidas por perfil na validação das demais lojas, com meta 0,70. Cada painel contém 145 escolhas; os números indicam séries. A primeira origem diferencia grupos, enquanto a segunda recomenda sSCycle em todos eles. Nomes abreviados dos perfis correspondem às regras operacionais descritas no método.}
\label{fig:profile_choices}
\end{figure}

A Figura~\ref{fig:profile_choices} localiza a diferenciação que produz o serviço adicional: na primeira origem, a regra por perfil substitui GA por sSCycle ou EOQCycle em parte das séries. Na segunda, identificar perfis não muda a política recomendada. Assim, a existência de grupos distintos não implica valor incremental da seleção em todo período.
\FloatBarrier
